% =========================================================================================
% Chapter 1 -- Introduction
% =========================================================================================
\section{Introduction}
\label{sec:introduction}

\subsection{Context: Emergency Medicine and POCUS}
\label{subsec:context}

In emergency medicine, clinical decisions are made in an environment where time is
limited, information is incomplete, and uncertainty is unavoidable. A patient arrives with
a few symptoms and abnormal vital signs, and the physician must rapidly determine whether
the situation is benign, potentially dangerous, or immediately life-threatening. In such
circumstances, every additional source of reliable information can influence both the
quality and the speed of the decision.

Among the tools available at the bedside, point-of-care ultrasound (\gls{pocus}) has
become an increasingly central component of emergency assessment. Unlike conventional imaging workflows, which may
require transporting the patient, scheduling specialised equipment and waiting for a
formal interpretation, \gls{pocus} allows the clinician to acquire images directly at the
bedside and use them within the same examination. It is portable, involves no ionising
radiation, costs comparatively little, and can provide information about the lungs, the
heart and the abdomen within minutes.

For the presentations that dominate emergency practice --- acute dyspnoea, chest pain,
abdominal pain, undifferentiated shock --- \gls{pocus} frequently answers the decisive
question faster than any other available investigation. Structured protocols have
formalised this role: the BLUE protocol for acute respiratory failure
\cite{lichtenstein2008blue} and the Focused Assessment with Sonography for Trauma
(\gls{fast}) in injury both reduce a complex diagnostic question to a short sequence of
bedside views, each with a defined interpretation.

\subsection{Problem Statement}
\label{subsec:problem}

\begin{keystatement}
The availability of an image does not, by itself, guarantee an accurate diagnosis.
\end{keystatement}

\gls{pocus} remains strongly dependent on the operator. Acquisition varies with probe
position, anatomical orientation, patient habitus, machine settings and technique.
Interpretation requires clinical expertise acquired over years of supervised practice, and
it is hardest precisely where it matters most: in emergency situations, where several
distinct conditions produce similar appearances, where the examination is performed under
time pressure, often at night, and often by a clinician whose primary specialty is not
radiology. The same clip that yields a confident diagnosis in experienced hands may yield
nothing in less experienced ones.

The challenge is therefore not to acquire more medical images, but to transform
heterogeneous observations into reliable and actionable clinical information. It grows
once the rest of the clinical picture is considered. An ultrasound examination is only one
part of the evidence available during an emergency evaluation. Vital signs, symptoms, age,
medical history and laboratory measurements provide additional evidence that may either
support or contradict what is seen on the screen. A patient may appear relatively stable
by their vital signs while the ultrasound reveals a potentially severe abnormality;
conversely, an isolated imaging finding may carry limited significance once the patient's
symptoms and physiological state are taken into account.

The problem this project addresses can therefore be stated as a single question:

\begin{keystatement}
How can multimodal artificial intelligence be integrated into point-of-care ultrasound to
improve diagnostic accuracy, assist clinicians at the bedside, and strengthen the quality
of emergency care?
\end{keystatement}

This question decomposes into four concrete difficulties, each of which shapes a design
decision later in this report.

\paragraph{Operator dependence.} Image acquisition and interpretation both vary with the
skill of the individual clinician. A system that assists interpretation must be robust to
variation in acquisition, not merely accurate on well-acquired images.

\paragraph{Interpretation under time pressure.} Emergency decisions are made in minutes.
An assistant whose output requires careful reading provides no benefit; its findings must
be immediately legible and must carry an explicit indication of how far they can be
trusted.

\paragraph{Fragmented information.} Vital signs, imaging, laboratory results and history
arrive separately and at different times. The clinical decision depends on their
combination, but each source is typically presented in isolation.

\paragraph{The cost of unwarranted confidence.} A decision-support system that presents an
uncertain finding with the same authority as a certain one is worse than no system at all,
because it introduces error while appearing to reduce it. This constraint proved to be the
single strongest influence on the architecture, and is the reason \gls{calibration},
evidence grading and the escalation policy are treated as core components rather than as
refinements.

\subsection{Motivation}
\label{subsec:motivation}

Recent advances in artificial intelligence (\gls{ai}) offer a way to approach this problem
from a multimodal perspective. Convolutional neural networks (\glspl{cnn}) can learn to
recognise patterns in ultrasound images and now match specialist performance on several
well-defined imaging tasks \cite{leclerc2019camus, isensee2021nnunet}. Supervised models
on structured data can analyse vital signs and patient characteristics. More recently,
large language models (\glspl{llm}) have demonstrated the ability to integrate
heterogeneous information and produce structured reasoning over complex textual and
numerical inputs \cite{chen2024huatuogpto1}.

Yet most published work stops at a single organ and a single question --- a model that
detects \gls{blines}, or one that segments the left ventricle --- evaluated in isolation.
A clinician facing an undifferentiated patient does not have a single question. They have
a patient, with vital signs, a history, laboratory results and an ultrasound probe, and
they must integrate all of it. This is the gap the project addresses: not the accuracy of
any individual model, but the absence of a layer that combines several of them with the
rest of the clinical record.

Placing a model in front of a physician does not, in itself, constitute a clinical
decision-support system. A useful emergency system must understand what it is looking at,
what it does not know, how reliable its information is, and when it should admit
uncertainty. A model that produces a confident answer from incomplete or contradictory
evidence is more dangerous than one that explicitly asks for additional information,
because it introduces error while appearing to remove it. This observation motivates the
central idea of the project:

\begin{keystatement}
Artificial intelligence should not replace the emergency clinician's reasoning. It should
act as a clinical co-pilot that organises evidence, highlights important findings,
identifies uncertainty, and supports safer decision-making.
\end{keystatement}

\subsection{Project Objectives}
\label{subsec:objectives}

\subsubsection{Business Objective}

The \gls{bo} is to improve the quality and speed of emergency decision-making for
critically ill patients. It decomposes into four measurable aims:

\begin{itemize}[label=$\bullet$]
  \setlength\itemsep{0.3em}
  \item \textbf{Reduce time-to-diagnosis.} An interpretation available during the
    examination itself, rather than after a formal report, shortens the interval between
    the patient arriving and the decisive question being answered. This is the aim most
    directly served by point-of-care imaging, and the one an assistant can most plausibly
    improve.
  \item \textbf{Reduce diagnostic error, particularly for less experienced operators.}
    Interpretation quality varies with training. An assistant that recognises defined
    pathologies narrows the gap between the confident reading an expert obtains and the
    uncertain one a junior clinician obtains from the same images.
  \item \textbf{Make uncertainty visible rather than implicit.} The value of the system
    lies as much in flagging what it does not know as in what it detects. A finding
    presented with an honest confidence, and an explicit statement of what has not been
    measured, allows the clinician to weigh it correctly instead of accepting or
    discarding it wholesale.
  \item \textbf{Preserve the clinician's authority over the decision.} The system produces
    a reasoned differential and, for ambiguous or high-stakes cases, exposes the reasoning
    so the physician can examine it before committing. It recommends next steps; it does
    not prescribe treatment and does not act autonomously.
\end{itemize}

\subsubsection{Data Science Objective}

The \gls{dso} is to deliver the technical components that make the above possible:

\begin{itemize}[label=$\bullet$]
  \setlength\itemsep{-0.2em}
  \item computer vision models that detect defined pathologies from \gls{pocus}, each
    returning a calibrated confidence rather than a bare score;
  \item an explainability layer relating model attention to anatomical regions a clinician
    can verify;
  \item a structured clinical state that records not only what was found, but what was not
    measured and what the models cannot exclude;
  \item a reasoning layer that converts imaging findings and clinical data into a ranked
    differential diagnosis with grounded justification;
  \item an escalation policy determining, from confidence, severity and data completeness,
    whether a case can be answered directly or requires further examination.
\end{itemize}

These reduce to three verbs that recur throughout the report:

\begin{enumerate}
  \setlength\itemsep{-0.2em}
  \item \textbf{Perceive} --- extract meaningful information from patient data and from
    \gls{pocus} images.
  \item \textbf{Reason} --- combine these observations with clinical knowledge and
    contextual information.
  \item \textbf{Escalate when necessary} --- recognise when the available evidence is
    insufficient or contradictory, rather than generating an unjustified conclusion.
\end{enumerate}

\begin{keystatement}
The objective is not to make artificial intelligence decide instead of the clinician, but
to make the available clinical evidence easier to understand, cross-check and act upon.
\end{keystatement}

\subsection{Proposed Approach}
\label{subsec:approach}

The project explores this concept through a modular multi-agent architecture. Rather than
developing a single model responsible for the entire diagnostic process, the system
separates the problem into specialised components, each answering a question it has the
information to answer. The overall organisation is shown in Fig.~\ref{fig:architecture};
each component is described in detail in Section~\ref{subsec:architecture}.

\begin{figure}[H]
  \centering
  \includegraphics[width=0.70\textwidth]{fig_architecture.png}
  \caption{Overall architecture of POCUS-Emergency. Perception, representation, reasoning
    and decision support are separated into distinct layers. The Clinical State Builder
    acts as a safety boundary between the perception models and the language model.}
  \label{fig:architecture}
\end{figure}

Two perception agents operate in parallel. The \textbf{Triage Agent} focuses on the
patient rather than the image, estimating an urgency level from the chief complaint, vital
signs and age. The \textbf{Ultrasound Agent} focuses exclusively on \gls{pocus} data,
routing an examination to the appropriate organ-specific model and returning structured
findings with confidence information and Gradient-weighted Class Activation Mapping
(\gls{gradcam}) explanations. Its objective is not
to decide whether a patient is urgent, but to answer a more fundamental question: what
does the ultrasound show?

Their outputs are combined by the \textbf{Clinical Reasoning Agent}, which weighs them
against laboratory results and history while accounting for missing information, model
limitations and contradictions between sources. A retrieval-augmented generation
(\gls{rag}) mechanism grounds the reasoning in a curated medical knowledge base
\cite{lewis2020rag}, and a locally deployed
medical language model transforms the structured evidence into a ranked differential
diagnosis with supporting arguments, explicit uncertainty and an escalation decision.

\paragraph{Reasoning under uncertainty.} A further objective is to address uncertainty
explicitly rather than implicitly. In real emergency medicine, the absence of information
is not equivalent to a normal result. A laboratory test that was never performed cannot be
interpreted as negative, and a model that was never trained to detect a particular
pathology cannot be used to exclude it. The system therefore incorporates an intermediate
Clinical State Builder that standardises the outputs of the different agents and
represents explicitly the information that is available and the information that is
missing, the confidence attached to each finding and whether that confidence is
calibrated, the stated limitations of each model, and any conflict between clinical and
imaging evidence. This representation acts as a safety boundary: it prevents the reasoning
layer from silently treating missing evidence as negative evidence, or from interpreting a
model outside the scope it was trained for.

\subsection{Report Organization}
\label{subsec:organization}

The report follows the order in which the system was built, so that each chapter depends
only on what precedes it.

\begin{itemize}[label=$\bullet$]
  \setlength\itemsep{-0.1em}
  \item \textbf{Chapter~\ref{sec:environment} --- Project Environment and Presentation.}
    The host organisation, the context of the internship, the system architecture in
    detail, and the development methodology.
  \item \textbf{Chapter~\ref{sec:sota} --- Preliminary Study and State of the Art.}
    Existing work in automated ultrasound interpretation and multimodal decision support,
    its limitations, and the positioning of this project against it.
  \item \textbf{Chapter~\ref{sec:data} --- Datasets and Data Preparation.} The datasets
    used, their characteristics, and the determination of the statistical unit that
    governs every result reported afterwards.
  \item \textbf{Chapter~\ref{sec:perception} --- Perception Agents.} Development and
    evaluation of the triage, lung, cardiac and gallbladder modules, including the
    experiments that failed and the corrections they prompted.
  \item \textbf{Chapter~\ref{sec:reasoning} --- Clinical Reasoning Agent.} The clinical
    state representation, conflict detection, retrieval, the medical language model,
    output validation and the escalation policy.
  \item \textbf{Chapter~\ref{sec:integration} --- System Integration and Prototype.} The
    end-to-end pipeline, the contract through which the agents communicate, and worked
    clinical scenarios.
  \item \textbf{Chapter~\ref{sec:results} --- Results and Discussion.} The measured
    results for each module, what they mean clinically, and the limitations that remain.
  \item \textbf{Chapter~\ref{sec:conclusion} --- Conclusion and Perspectives.} The
    contributions, the objectives achieved, and directions for future work.
\end{itemize}
